\documentclass[11pt]{article}

\usepackage{amsmath, amssymb, amsthm}
\usepackage{geometry}
\usepackage{graphicx}
\usepackage{bm}

\geometry{margin=1in}

\title{\textbf{CSE400 --- Lecture 15--17 Scribe}}
\author{Instructor: Prof. Dhaval Patel, PhD}
\date{March 10, 2026}

\begin{document}

\maketitle

\noindent \textbf{Topic:} Joint Random Variables; Joint Distributions; Joint CDF

\vspace{0.5cm}

%--------------------------------------------------

\section{Overview and Objective}

This lecture studies:

\begin{enumerate}
    \item Joint random variables
    \item Joint probability mass function (PMF)
    \item Joint probability density function (PDF)
    \item Joint cumulative distribution function (CDF)
    \item Worked examples (discrete and continuous)
\end{enumerate}

\textbf{Key Objective:}  
Given multiple random variables, analyze their \textbf{combined behavior} and compute probabilities involving them.

\textbf{Assumption:}  
All random variables are defined on a common probability space.

\vspace{0.5cm}

%--------------------------------------------------

\section{Why Study Joint Random Variables?}

In real-world applications, we are rarely interested in a single random variable.

Instead, we analyze relationships between multiple variables.

\vspace{0.3cm}

\subsection{Example 1: Smart Transportation}

\[
X = \text{Vehicle Speed}, \quad Y = \text{Traffic Density}
\]

\textbf{Question:} Can speed be high when density is high?

\textbf{Observation:}  
The variables are dependent (correlated).

\vspace{0.2cm}

\subsection{Example 2: Wireless Networks}

\[
X = \text{Signal-to-Noise Ratio (SNR)}, \quad Y = \text{Throughput}
\]

\textbf{Observation:}
\begin{itemize}
    \item Throughput depends on SNR
    \item External factors (congestion) affect dependency
\end{itemize}

\vspace{0.2cm}

\subsection{Example 3: Weather System}

\[
X = \text{Rainfall}, \quad Y = \text{Temperature}
\]

\vspace{0.2cm}

\subsection{Example 4: Drone Delivery}

\[
X = \text{Wind Speed}, \quad Y = \text{Delivery Time}
\]

\vspace{0.2cm}

\subsection{Example 5: Rolling Two Dice}

\[
X = \text{First die outcome}, \quad Y = \text{Second die outcome}
\]

This requires computation of joint probabilities such as:

\[
P(X + Y = 7), \quad P(X > Y)
\]

\vspace{0.5cm}

%--------------------------------------------------

\section{Discrete Case: Joint PMF}

\subsection{Definition}

Let $X$ and $Y$ be discrete random variables.

The \textbf{joint PMF} is defined as:

\[
p_{X,Y}(x,y) = P(X = x, Y = y)
\]

\vspace{0.3cm}

\subsection{Properties}

\begin{enumerate}
    \item Non-negativity:
    \[
    p_{X,Y}(x,y) \ge 0
    \]
    
    \item Total probability:
    \[
    \sum_x \sum_y p_{X,Y}(x,y) = 1
    \]
\end{enumerate}

\vspace{0.3cm}

\subsection{Joint PMF Table Representation}

\begin{itemize}
    \item Rows correspond to values of $X$
    \item Columns correspond to values of $Y$
    \item Each cell contains $P(X=x, Y=y)$
\end{itemize}

\vspace{0.5cm}

%--------------------------------------------------

\section{Worked Example 1: Rolling Two Dice}

\subsection{Problem Setup}

Two fair dice are rolled.

\[
X = \text{value of first die}, \quad Y = \text{value of second die}
\]

\vspace{0.3cm}

\subsection{Step 1: Sample Space}

The sample space consists of ordered pairs:

\[
(x,y), \quad x,y \in \{1,2,3,4,5,6\}
\]

Total outcomes:

\[
6 \times 6 = 36
\]

\vspace{0.3cm}

\subsection{Step 2: Joint PMF}

Since dice are fair:

\[
P(X=x, Y=y) = \frac{1}{36}
\]

for all valid $(x,y)$.

\vspace{0.3cm}

\subsection{Step 3: Event $X + Y = 7$}

List all valid outcomes:

\[
(1,6), (2,5), (3,4), (4,3), (5,2), (6,1)
\]

Number of outcomes = 6

\[
P(X+Y=7) = \frac{6}{36} = \frac{1}{6}
\]

\vspace{0.3cm}

\subsection{Step 4: Event $X > Y$}

Count all pairs such that $x > y$:

\[
(2,1), (3,1), (3,2), (4,1), (4,2), (4,3),
\]
\[
(5,1), (5,2), (5,3), (5,4),
\]
\[
(6,1), (6,2), (6,3), (6,4), (6,5)
\]

Total = 15

\[
P(X > Y) = \frac{15}{36}
\]

\vspace{0.5cm}

%--------------------------------------------------

\section{Continuous Case: Joint PDF}

\subsection{Definition}

For continuous random variables:

\[
f_{X,Y}(x,y)
\]

is the joint probability density function.

\vspace{0.3cm}

\subsection{Properties}

\begin{enumerate}
    \item Non-negativity:
    \[
    f_{X,Y}(x,y) \ge 0
    \]
    
    \item Total probability:
    \[
    \iint_{-\infty}^{\infty} f_{X,Y}(x,y)\,dx\,dy = 1
    \]
\end{enumerate}

\vspace{0.3cm}

\subsection{Geometric Interpretation}

Probability over a region $R$:

\[
P((X,Y)\in R) = \iint_R f_{X,Y}(x,y)\,dx\,dy
\]

This corresponds to the \textbf{volume under the surface}.

\vspace{0.5cm}

%--------------------------------------------------

\section{Worked Example 2: Continuous Case}

\subsection{Step 1: Identify Region}

Determine limits of integration for $x$ and $y$.

\vspace{0.3cm}

\subsection{Step 2: Verify PDF}

Check:

\[
\iint f_{X,Y}(x,y)\,dx\,dy = 1
\]

\vspace{0.3cm}

\subsection{Step 3: Compute Probability}

\[
P((X,Y)\in R) = \iint_R f_{X,Y}(x,y)\,dx\,dy
\]

\vspace{0.3cm}

\subsection{Step 4: Evaluate Integral}

Perform integration step-by-step using limits.

\vspace{0.5cm}

%--------------------------------------------------

\section{Joint Cumulative Distribution Function (CDF)}

\subsection{Definition}

\[
F_{X,Y}(x,y) = P(X \le x, Y \le y)
\]

\vspace{0.5cm}

%--------------------------------------------------

\section{Joint CDF: Continuous Case}

\[
F_{X,Y}(x,y) =
\int_{-\infty}^{x}
\int_{-\infty}^{y}
f_{X,Y}(u,v)\,dv\,du
\]

\vspace{0.5cm}

%--------------------------------------------------

\section{Joint CDF: Discrete Case}

\[
F_{X,Y}(x,y) =
\sum_{u \le x}
\sum_{v \le y}
P(X=u, Y=v)
\]

\vspace{0.5cm}

%--------------------------------------------------

\section{Properties of Joint CDF}

\begin{enumerate}
    \item $0 \le F_{X,Y}(x,y) \le 1$
    
    \item Monotonicity:
    \[
    F(x_1,y_1) \le F(x_2,y_2)
    \quad \text{if } x_1 \le x_2, y_1 \le y_2
    \]
    
    \item Limits:
    \[
    \lim_{x,y \to \infty} F_{X,Y}(x,y) = 1
    \]
    
    \item Boundary behavior:
    \[
    \lim_{x \to -\infty \text{ or } y \to -\infty} F_{X,Y}(x,y) = 0
    \]
\end{enumerate}

\vspace{0.5cm}

%--------------------------------------------------

\section{Logical Flow Summary}

\begin{enumerate}
    \item Define joint random variables
    \item Construct joint PMF/PDF
    \item Compute probabilities using sums/integrals
    \item Define joint CDF
    \item Use CDF for cumulative probabilities
\end{enumerate}

\vspace{0.5cm}

%--------------------------------------------------

\section{Key Results for Exams}

\begin{enumerate}
    \item Joint PMF:
    \[
    P(X=x, Y=y)
    \]
    
    \item Joint PDF:
    \[
    P((X,Y)\in R) = \iint_R f_{X,Y}(x,y)\,dx\,dy
    \]
    
    \item Joint CDF:
    \[
    F_{X,Y}(x,y) = P(X \le x, Y \le y)
    \]
\end{enumerate}

\end{document}